\section{Base Values\label{sec:BaseValues}}
\hypertarget{def-basevalueterm}{}
Each type, with the exceptions stated below, has a \basevalueterm,
which is used to initialize storage elements (either local or global),
if an initializer is not supplied.

\RequirementDef{NoBaseValue}
The following types do not have a \basevalueterm{}:
\begin{itemize}
    \item \parameterizedintegertypesterm{} ---
          only subprogram parameters can be parameterized integers, and since they will be initialized by their
          invocation, there is no need to define a \basevalueterm{} for them;
    \item a \wellconstrainedintegertypeterm{} whose list of constraints
          represents the empty set;
    \item a \bitvectortypeterm{} whose length is negative.
\end{itemize}

\ExampleDef{Base Values}
\listingref{base-values} shows a specification with examples of well-typed \basevalueterm{}
for various types, followed by the output to the console.
\ASLListing{Well-typed Base Values}{base-values}{\typingtests/TypingRule.BaseValue.asl}
% CONSOLE_BEGIN aslref \typingtests/TypingRule.BaseValue.asl
\begin{Verbatim}[fontsize=\footnotesize, frame=single]
global_base = 0, unconstrained_integer_base = 0, constrained_integer_base = -3
bool_base = FALSE, real_base = 0, string_base = , enumeration_base = RED
bits_base = 0x00
tuple_base = (0, -3, RED)
record_base      = {data=0x00, time=0, flag=FALSE}
record_base_init = {data=0x00, time=0, flag=FALSE}
exception_base = {msg=}
integer_array_base = [[0, 0, 0, 0]]
enumeration_array_base = [[RED=0, GREEN=0, BLUE=0]]
\end{Verbatim}
% CONSOLE_END

\listingref{base-values-parameterised} shows a specification that relies on the base value of a \bitvectortypeterm{} whose width is a \parameterizedintegertypeterm{}.
\ASLListing{Base Value for Parameterized Bitvector Width}{base-values-parameterised}{\typingtests/TypingRule.BaseValue.parameterized.asl}

\ExampleDef{Types Without Base Value}
\listingref{base-values-bad-negative-width} shows an ill-typed specification
where the width of a bitvector is negative.
\ASLListing{No Base Value for Bitvectors of Negative Width}{base-values-bad-negative-width}{\typingtests/TypingRule.BaseValue.bad_negative_width.asl}

\listingref{base-values-bad-empty-type} shows an ill-typed specification
where the constraint \verb|5..0| represents an empty set.
Therefore, the domain of values for the type \verb|integer{5..0}| is empty,
which negates the possibility of having a \basevalueterm.
\ASLListing{No Base Value for an Empty Integer Type}{base-values-bad-empty-type}{\typingtests/TypingRule.BaseValue.bad_empty.asl}

\RenderRelation{base_value}

\TypingRuleDef{BaseValue}
See \ExampleRef{Base Values} and \ExampleRef{Types Without Base Value}.


\RenderProseAndFormally{base_value}

\TypingRuleDef{ConstraintAbsMin}
\RenderRelation{constraint_abs_min}

\ExampleDef{Minimal Absolute Value in a Constraint List}
The minimal absolute value of \verb|{7, -2}| is \verb|-2|.\\
%
The minimal absolute value of \verb|{2, -2}| is \verb|2|.\\
%
The minimal absolute value of \verb|{-2..2, 5}| is \verb|0|.


\RenderProseAndFormally{constraint_abs_min}

\TypingRuleDef{ListMinAbs}
\RenderRelation{list_min_abs}

\ExampleDef{Minimal Absolute Value}
The minimal absolute value of $[9, -3]$ is $-3$,
and the minimal absolute value of $[2, -2]$ is $2$.


\RenderProseAndFormally{list_min_abs}
